\begin{pastexam}[label={px:ncku-110-prog-08}]{成大資訊聯招 110 學年度 第 8 題}
（填空）

(10\%) Given three strings $x[0,\ldots,n-1]$, $y[0,\ldots,m-1]$ and $z[0,\ldots,r-1]$. We say
that $z$ is a shuffle of $x$ and $y$ if it contains all characters of $x$ and $y$ and the
left-to-right ordering of the characters from $x$ and the characters from $y$ is preserved. For
example, ``NcCKsiUe'' is a shuffle of ``NCKU'' and ``csie''. The dynamic programming algorithm
uses a table $S$ to check $x$, $y$ and $z$, where $S[i][j]$ is true if and only if the first $i+j$
characters of $z$ are a shuffle of the first $i$ characters of $x$ together with the first $j$
characters of $y$. Complete the following pseudocode for checks whether $z$ is a shuffle of $x$
and $y$ by filling (1), (2) and (3). (please use C-style expression)

\begin{tabbing}
\quad\=\quad\=\quad\=\quad\=\kill
isShuffle$(x,y,z)$\\
\>Let $S[0\ldots n][0\ldots m]$ be a new table\\
\>$S[0][0]=$ true\\
\>if $r\neq n+m$\\
\>\>\textbf{return} false\\
\>for $i=1$ to $n$\\
\>\>$S[i][0]=$ \blank{1}~(3\%)\\
\>for $j=1$ to $m$\\
\>\>$S[0][j]=$ \blank{2}~(3\%)\\
\>for $i=1$ to $n$\\
\>\>for $j=1$ to $m$\\
\>\>\>$S[i][j]=$ \blank{3}~(4\%)\\
\>\textbf{return} $S[n][m]$
\end{tabbing}
\end{pastexam}

\begin{pastsoln}
答案（C 的寫法；字串從 $0$ 開始編號）：
\begin{itemize}
\item[(1)] \texttt{S[i-1][0] \&\& x[i-1] == z[i-1]}
\item[(2)] \texttt{S[0][j-1] \&\& y[j-1] == z[j-1]}
\item[(3)] \texttt{(S[i-1][j] \&\& x[i-1] == z[i+j-1]) || (S[i][j-1] \&\& y[j-1] == z[i+j-1])}
\end{itemize}

\textbf{子問題}（資結書第~18.1.5~節的步驟）。 $S[i][j]$：$z$ 的前 $i + j$ 個字元能不能由 $x$ 的
前 $i$ 個和 $y$ 的前 $j$ 個交錯而成。 $z$ 的前 $i + j$ 個字元裡，最後一個是 $z[i+j-1]$（從 $0$
開始編號）。 它不是來自 $x$ 就是來自 $y$：
\begin{itemize}
\item 來自 $x$：它必須是 $x$ 前 $i$ 個字元的最後一個 $x[i-1]$（順序要保持），而且去掉它以後，$z$
的前 $i + j - 1$ 個是 $x$ 前 $i - 1$ 個和 $y$ 前 $j$ 個的交錯，也就是 $S[i-1][j]$。
\item 來自 $y$：同理，要 $y[j-1] = z[i+j-1]$ 且 $S[i][j-1]$。
\end{itemize}
兩種情形任一成立即可，用 \texttt{||} 連起來，就是 (3)。 兩種情形都可能成立（$x[i-1]$ 和
$y[j-1]$ 是同一個字元時），所以不能只檢查一邊。

\textbf{邊界}。 $j = 0$ 時 $y$ 一個字元都沒用，$z$ 的前 $i$ 個必須逐字等於 $x$ 的前 $i$ 個：前
$i - 1$ 個要相等（$S[i-1][0]$），第 $i$ 個也要相等（$x[i-1] = z[i-1]$），就是 (1)。 $i = 0$ 時
對稱，就是 (2)。 $S[0][0] = $ true 是兩個空字串交錯成空字串。 填表之前先檢查長度：$r \ne n + m$
時不可能是交錯，直接回傳 false。

\textbf{時間}。 表有 $(n+1)(m+1)$ 格，每格 $O(1)$，總共 $O(nm)$。

\textbf{用題目的例子檢查}。 $x = $ NCKU、$y = $ csie、$z = $ NcCKsiUe，$n = m = 4$、$r = 8$。 表
$S$（列 $i = 0..4$、行 $j = 0..4$，T 是 true，空白是 false）：
\begin{center}
\small
\begin{tabular}{c|ccccc}
$S$ & $j{=}0$ & 1 & 2 & 3 & 4 \\ \hline
$i{=}0$ & T & & & & \\
1 & T & T & & & \\
2 & & T & & & \\
3 & & T & T & T & \\
4 & & & & T & T \\
\end{tabular}
\end{center}
例如 $S[1][1]$：$z[1] = $ c，從 $S[1][0] = $ T 且 $y[0] = $ c $= z[1]$ 得 T。 $S[3][3]$：$z[5] = $ i，
從 $S[3][2] = $ T 且 $y[2] = $ i 得 T。 $S[4][4]$：$z[7] = $ e $= y[3]$，而 $S[4][3] = $ T，得 T。
所以 NcCKsiUe 是 NCKU 和 csie 的交錯，和題目說的一致。 走回去的路徑是
$(0,0) \to (1,0) \to (1,1) \to (2,1) \to (3,1) \to (3,2) \to (3,3) \to (4,3) \to (4,4)$，
對應 $z$ 的字元依序取自 $x, y, x, x, y, y, x, y$：N、c、C、K、s、i、U、e。
\end{pastsoln}
